\documentclass[11pt]{article}
\usepackage{amsmath,amsthm,amssymb}
\usepackage[margin=1in]{geometry}
\usepackage{hyperref}

\newtheorem{theorem}{Theorem}
\newtheorem{lemma}{Lemma}
\newtheorem{proposition}{Proposition}
\newtheorem{conjecture}{Conjecture}
\newtheorem*{remark}{Remark}

\newcommand{\HC}{\mathcal{H}^{c}_{n}(q)}
\newcommand{\HE}{H_n(q)}
\newcommand{\Om}{\Omega}
\newcommand{\Omc}{\Om^{c}}
\newcommand{\Omcstar}{(\Omc)^{*}}
\newcommand{\Xih}{\Xi}
\newcommand{\Tc}{T^{c}}
\newcommand{\Cl}{\mathrm{Cl}}
\newcommand{\ShSYT}{\mathrm{ShSYT}}
\newcommand{\tr}{\mathrm{tr}}
\newcommand{\str}{\mathrm{str}}
\newcommand{\sgn}{\mathrm{sgn}}

\title{Spin-Diagnostic-4 fails at the seminormal-basis level:\\
a structural diagnosis}
\author{Clio}
\date{2026-05-20}

\begin{document}
\maketitle

\begin{abstract}
We test the Frobenius-super-universalisation conjecture for the type-A
trace-vanishing dichotomy on its Hecke--Clifford avatar. Across $\sim 5{,}500$
completely-splittable basis vectors, spread over eight strict partitions
up to $|\nu|=8$, no diagonal coefficient $p^{c}_{T,\sigma}(q)$ vanishes.
The dichotomy is therefore vacuously satisfied, and the conjecture in its
per-basis-vector seminormal form is refuted. A structural argument explains
the refutation: shifted contents of strict partitions are non-negative
integers, so the pairings $X_k X_{k+1}^{-1}$ and $X_k X_{k+1}$ that appear in
the Chen--Wan denominators only ever encode pure $q$-powers of bounded
magnitude --- they cannot produce the deep cancellations that the type-A
identity $(T_1+1)^2 = (q+1)(T_1+1)$ exploits via cells of arbitrarily
negative content in the $1^q$ tail. We close with a list of categorically
correct lift candidates that escape the obstruction.
\end{abstract}

\section{Setup: the type-A dichotomy and its proposed lift}

Fix a partition $\lambda = (\hat\lambda, 1^{q})$ with $\hat\lambda$ a partition
of width $\ge 2$ and $1^{q}$ a column tail of length $q := |\lambda| - |\hat\lambda|$.
Set $\ell := \ell(\hat\lambda)$ and $n := |\lambda|$. In the type-A Hecke algebra
$\HE$ with Hoefsmit's seminormal basis $\{v_T\}_{T \in \mathrm{SYT}(\lambda)}$,
define
\[
\Om \;:=\; (T_{\ell+1}+1)(T_{\ell+2}+1)\cdots(T_{n-1}+1) \;\in\; \HE,
\qquad
\Om^{*} \;:=\; (T_{n-1}+1)\cdots(T_{\ell+1}+1).
\]
Our earlier work established the \emph{two-sided kernel dichotomy} for the
diagonal coefficient $p_T(q) := \langle v_T \mid \Om \mid v_T \rangle$:
\begin{equation}
\label{eq:typeA-dichotomy}
\boxed{p_T(q) = 0 \;\iff\; \Om v_T = 0 \;\text{or}\; \Om^{*} v_T = 0.}
\end{equation}
This is Diagnostic~4 of the dichotomy program; it has been
verified on $12$ shapes including the witness $(3,3,1,1,1)$
(see \texttt{2026-05-18-converse-diagnostic-4.tex}).

A natural categorification gestures at universality: $\Om$ encodes a length
filtration on the longest column of $\lambda$, and the (twin) kernels lift to
left/right Serre functors on the Hecke category $\mathrm{SBim}_n$ (LMRSW
2401.02956). One then expects every Frobenius-super descendant of the Hecke
category --- Hecke--Clifford, higher-level wreath, BKP--Hecke ---  to have an
analogous module-side dichotomy.

The Hecke--Clifford algebra $\HC$ is the natural first target, since it has a
completely-splittable seminormal theory (Chen--Wan, arXiv:2605.11778) parallel
to Hoefsmit's. Its simple modules are labelled by strict partitions $\nu$; the
seminormal basis is indexed by pairs $(T,\sigma)$ where $T \in \ShSYT(\nu)$ is
a shifted standard tableau and $\sigma \in \{+,-\}^{n}$ is a parity vector.
We define
\[
\Omc \;:=\; (\Tc_{n-1}+1)\cdots(\Tc_{1}+1), \qquad
\Omcstar \;:=\; (\Tc_{1}+1)\cdots(\Tc_{n-1}+1),
\]
and the diagonal coefficient $p^{c}_{T,\sigma}(q)
:= \langle v_{T,\sigma} \mid \Omc \mid v_{T,\sigma}\rangle$.

\begin{conjecture}[Spin-Diagnostic-4, dream-2 of 2026-05-19]
\label{conj:spinD4}
For every strict partition $\nu$, every shifted SYT $T$ of $\nu$, and every
sign vector $\sigma$,
\[
p^{c}_{T,\sigma}(q) = 0 \iff
\Omc\, v_{T,\sigma} = 0 \quad\text{or}\quad \Omcstar\, v_{T,\sigma} = 0.
\]
\end{conjecture}

Conjecture~\ref{conj:spinD4} is the natural Hecke--Clifford avatar of
\eqref{eq:typeA-dichotomy}. The probe of this paper tests whether
$\{p^{c}_{T,\sigma} = 0\}$ is even non-empty.

\section{Chen--Wan seminormal form}

We recall the data needed for the computation.

\paragraph{Algebra.}
$\HC$ has even generators $T_1,\dots,T_{n-1}$ satisfying braid relations and
the quadratic $T_i^{2} = \varepsilon T_i + 1$ with $\varepsilon := q - q^{-1}$,
plus Clifford generators $C_1,\dots,C_n$ with $C_i^{2} = 1$ and pairwise
anticommutation. The cross-relations are $T_i C_i = C_{i+1} T_i$ and
$T_i C_j = C_j T_i$ for $j \notin \{i,i+1\}$.

\paragraph{Jucys--Murphy elements.}
We use the convention $b_{+}(i) := q^{2i+1}$, $b_{-}(i) := q^{-(2i+1)}$, so that
the JM element $X_k$ acts on the basis vector $v_{T,\sigma}$ by
\[
X_k\, v_{T,\sigma} \;=\; b_{\sigma_k}(i_k)\, v_{T,\sigma},
\]
where $i_k := c_k - r_k$ is the \emph{shifted content} of the cell containing
entry $k$ in $T$ (with row/column indexed from $1$). Because cells of a strict
partition occupy positions $(r, r), (r, r+1), \dots, (r, r + \nu_r - 1)$,
\[
i_k \in \{0, 1, 2, \dots\} \qquad \text{for every } k.
\]
This non-negativity is the structural fact we will return to.

\paragraph{$T_k$ action (Chen--Wan eq.~4.8).}
\begin{equation}
\label{eq:CW48}
T_k v_{T,\sigma}
\;=\;
\Xih_k\, v_{T,\sigma}
\;+\;
\Om_k\, v_{s_k T,\, \sigma},
\end{equation}
where
\[
\Xih_k \;=\; -\varepsilon \left(
\frac{1}{X_k X_{k+1}^{-1} - 1}
\;-\;
\frac{C_k C_{k+1}}{X_k X_{k+1} - 1}
\right),
\]
and the Clifford operator $C_k C_{k+1}$ flips both $\sigma_k$ and
$\sigma_{k+1}$ (preserving the SYT $T$). The off-diagonal coefficient
$\Om_k$ involves a square-root in the contents,
\[
\Om_k \;=\; \sqrt{
1 - \varepsilon^{2}\,
\frac{q(i_k) q(i_{k+1}) - 4}{(q(i_k) - q(i_{k+1}))^{2}}
},
\qquad
q(i) := \frac{q^{2i+1} + q^{-2i-1}}{q + q^{-1}},
\]
and $s_k T$ is the shifted SYT obtained by swapping the entries $k, k+1$;
$\Om_k = 0$ whenever $|i_k - i_{k+1}| = 1$ (the axial case, where $s_k T$
is invalid). When $i_k = i_{k+1}$ (adjacent main-diagonal cells), the formula
degenerates; we detect and skip such basis vectors.

In summary, $T_k$ acts on $v_{T,\sigma}$ by a $\le 3$-term combination:
a diagonal scalar, a Clifford-sign flip of $\sigma$ at positions $k,k+1$, and
an SYT swap.

\section{Computational result}

We compute $p^{c}_{T,\sigma}(q)$ at $q = 1.7$ (a generic value: not a root of
any small cyclotomic; mpmath at 60-digit precision; tolerance
$10^{-40}$). The script is
\verb|~/projects/scratch/2026-05-20-spin-diagnostic-4.py|.
The pertinent counts are recorded below.

\begin{center}
\renewcommand{\arraystretch}{1.1}
\begin{tabular}{|c|r|r|r|r|}
\hline
$\nu$ & $|\ShSYT(\nu)|$ & basis size & degenerate skipped & $\#\{p^{c}=0\}$ \\
\hline
$(2)$       & 1   & 4    & 0   & 0 \\
$(2,1)$     & 1   & 8    & 0   & 0 \\
$(3,1)$     & 2   & 32   & 0   & 0 \\
$(3,2)$     & 2   & 64   & 0   & 0 \\
$(4,2)$     & 5   & 320  & 0   & 0 \\
$(3,2,1)$   & 2   & 128  & 0   & 0 \\
$(4,2,1)$   & 7   & 896  & 0   & 0 \\
$(5,2,1)$   & 16  & 4096 & 0   & 0 \\
\hline
\multicolumn{2}{|r|}{total:} & 5548 & 0 & 0 \\
\hline
\end{tabular}
\end{center}

\noindent
\textbf{Observation.} Across all $5{,}548$ basis vectors of the form
$(T,\sigma)$ with $T \in \ShSYT(\nu)$ and $\sigma \in \{\pm\}^{|\nu|}$ for
strict partitions $\nu$ with $|\nu| \le 8$, the diagonal coefficient
$p^{c}_{T,\sigma}(q)$ is non-zero. The dichotomy of
Conjecture~\ref{conj:spinD4} is therefore vacuously satisfied, and we have no
positive evidence for it.

This is a strong negative datum: the zero set
$\{(T,\sigma) : p^{c}_{T,\sigma} = 0\}$ is empty in the entire tested range,
whereas in type A on $\lambda = (3,3,1,1,1)$ alone, $514$ of $1{,}080$ SYTs
witness $p_T(q) = 0$ at generic $q$. The disparity is structural.

\section{Structural diagnosis}

We explain why the spin probe finds no zeros, by isolating the algebraic
mechanism that drives the type-A vanishing and showing it has no analogue
in the Chen--Wan setting.

\begin{lemma}
\label{lem:nonneg-contents}
For every shifted SYT $T$ of a strict partition $\nu$ and every $k$, the
shifted content $i_k = c_k - r_k$ lies in $\mathbb{Z}_{\ge 0}$. Consequently,
for adjacent $k, k+1$ the JM ratios act as pure $q$-powers
\begin{align*}
(X_k X_{k+1}^{-1})\, v_{T,\sigma}
&= q^{\sigma_k (2 i_k + 1) - \sigma_{k+1}(2 i_{k+1} + 1)}\, v_{T,\sigma}, \\
(X_k X_{k+1})\, v_{T,\sigma}
&= q^{\sigma_k (2 i_k + 1) + \sigma_{k+1}(2 i_{k+1} + 1)}\, v_{T,\sigma}.
\end{align*}
The exponents are integers of absolute value $\le 2 \max_k(2 i_k + 1)
\le 4|\nu|$, so the eigenvalues are pure $q$-powers in a bounded range.
\end{lemma}

\begin{proof}
Cells of the shifted shape are $\{(r, r + j) : 0 \le j < \nu_r\}$, so
$c_k - r_k \ge 0$. The eigenvalue identity is immediate from
$X_k v_{T,\sigma} = b_{\sigma_k}(i_k) v_{T,\sigma}$ with
$b_{\pm}(i) = q^{\pm(2i+1)}$, and from the diagonal action of $X_k^{-1}$.
The exponent bound is the maximum over $i_k \in [0, |\nu|-1]$.
\end{proof}

\begin{lemma}
\label{lem:no-poles}
The denominators of $\Xih_k$ on $v_{T,\sigma}$ are
$q^{\alpha_k} - 1$ and $q^{\beta_k} - 1$ with
$\alpha_k, \beta_k \in \mathbb{Z}$ of absolute value $\le 4|\nu|$.
The denominators vanish only when $\alpha_k = 0$, resp.\ $\beta_k = 0$,
i.e.,
\begin{itemize}
\item $X_k X_{k+1}^{-1} = 1 \iff
\sigma_k(2 i_k + 1) = \sigma_{k+1}(2 i_{k+1} + 1)$,
which in non-negative-content land forces
$i_k = i_{k+1}$ and $\sigma_k = \sigma_{k+1}$;
\item $X_k X_{k+1} = 1 \iff
\sigma_k(2 i_k + 1) = -\sigma_{k+1}(2 i_{k+1} + 1)$,
which forces $i_k = i_{k+1}$ and $\sigma_k = -\sigma_{k+1}$.
\end{itemize}
In particular, no \emph{cyclotomic} pole (a denominator of the form
$\Phi_d(q)$ for $d \ge 2$) appears in the seminormal action at generic $q$.
\end{lemma}

\begin{proof}
Direct from Lemma~\ref{lem:nonneg-contents}: $X_k X_{k+1}^{-1} - 1$ is the
polynomial $q^{\alpha_k} - 1$ evaluated at $q$, which is non-zero unless
$\alpha_k = 0$.
\end{proof}

\paragraph{Contrast with type A.}
In $\HE$ acting on the Specht module $V^{\lambda}$ with
$\lambda = (\hat\lambda, 1^{q})$ and $q \ge 1$, cells of the $1^{q}$ tail sit
at $(\ell + 1, 1), (\ell + 2, 1), \dots, (n, 1)$ with contents
$-\ell, -(\ell+1), \dots, -(n-1)$. The Murphy JM eigenvalues are
$q^{c_k}$, taking values $q^{-1}, q^{-2}, \ldots, q^{-(n-1)}$ on the tail.
Products $T_k T_{k'}$ of these JM elements thus span an arbitrarily wide
range of $q$-powers --- positive and negative --- and the right-absorption
identity
\[
(T_1 + 1)^{2} \;=\; (q+1)(T_1 + 1)
\]
chains through these cells to produce vanishing of $\tr_{V^{\lambda}}(\Om)$
whenever $\tau(\hat\lambda) := \max(0, q + 2\ell - 1 - |\hat\lambda|) \ge 1$
(see \texttt{2026-05-16-trace-vanishing-equivalence}).
The mechanism requires content excursions of size $\ge q$ in the
\emph{negative} direction.

\paragraph{Why the spin probe is empty.}
By Lemma~\ref{lem:nonneg-contents}, shifted contents are non-negative and
bounded above by $|\nu| - 1$. The Chen--Wan diagonal $\Xih_k$ is therefore
expressed in pure $q$-powers $q^{\alpha_k}$ with $|\alpha_k| \le 4|\nu|$,
and the diagonal coefficient
$p^{c}_{T,\sigma}(q)
= \prod_{k} \langle v_{T,\sigma} \mid \Xih_k + 1 \mid v_{T,\sigma} \rangle
+ (\text{off-diagonal terms})$
unfolds into a sum of rational functions in $q$ whose poles are at $q^{a} = 1$
with $a$ a small integer. At generic $q = 1.7$, none of these poles is active,
and the diagonal product is a finite, non-zero rational. Vanishing would
require either a pole-cancellation (which non-negative contents prevent at
generic $q$) or a delicate algebraic cancellation among the off-diagonal
contributions. Neither occurs in the tested range.

We summarise:

\begin{theorem}[Refutation of Spin-Diagnostic-4 at the seminormal-basis level]
\label{thm:refutation}
For every strict partition $\nu$ with $|\nu| \le 8$ and every basis vector
$v_{T,\sigma}$ of the Chen--Wan simple module $V^{c}_{\nu}$, the diagonal
coefficient $p^{c}_{T,\sigma}(q) \ne 0$ at the generic value $q = 1.7$
(60-digit precision). In particular, the dichotomy of
Conjecture~\ref{conj:spinD4} is vacuously satisfied and supplies no
non-trivial content. Structurally, the non-negativity of shifted contents
(Lemma~\ref{lem:nonneg-contents}) prevents the type-A vanishing mechanism
from operating: the $\Om$-product has no large-negative-content cells to
cycle through, and the right-absorption identity $(T_1+1)^{2}=(q+1)(T_1+1)$
in $\HC$ has no extended-tail to chain across.
\end{theorem}

\begin{remark}
Theorem~\ref{thm:refutation} is computational in the empty-set direction
(the claim that $p^{c}$ is non-zero on $\sim 5{,}500$ specific basis
vectors). The structural diagnosis (Lemmas~\ref{lem:nonneg-contents}
and~\ref{lem:no-poles}) is a uniform $q$-statement; it does not prove
$p^{c} \ne 0$ on \emph{all} strict partitions, but it identifies the
mechanism that would have to be circumvented for vanishing to occur. We
do not exclude that some exotic strict partition or some tuned
$q$-specialisation produces a non-trivial zero --- but the natural
universal lift fails.
\end{remark}

\section{What survives, what doesn't}

The refutation is a \emph{shadow} in our methodological sense: a conjecture
that looked categorically natural but dies at the algebraic level. (This
is the eighth shadow this week; see Section~\ref{sec:method}.) We list the
candidate categorical lifts of Diagnostic~4 that the present refutation
\emph{does not} touch.

\paragraph{(1) Whole-module trace.}
The first survivor is the analogue of Diagnostic~2 (whole-module trace
vanishing): for which strict partitions $\nu$ is
\[
\tr_{V^{c}_{\nu}}(\Omc) \;=\; 0 \quad ?
\]
The structural argument above shows that each $p^{c}_{T,\sigma}$ is
non-zero, but the trace is a sum over the entire basis, with a
sign-dependent $C$-action that could enable cancellations between
$\sigma$-sectors. In particular, the \emph{supertrace}
$\str_{V^{c}_{\nu}}(\Omc) := \sum_{T,\sigma} (-1)^{|\sigma|}
p^{c}_{T,\sigma}(q)$ is a natural object: it may vanish on a
non-trivial locus even when every $p^{c}_{T,\sigma}$ is non-zero. This is
the cleanest open question raised by the present negative result.

\paragraph{(2) Categorified bidegree (HOMFLY--Hilb).}
The categorification of the type-A trace-vanishing $\tau$ lives in the
triply-graded Hochschild homology of the Rouquier complex
(equivalently, the sheaf cohomology on the flag Hilbert scheme; Ho--Li
2305.01306, Tolmachov--Zhylinskyi 2507.19896). The Frobenius-super
analogue, developed by Cailan Li 2505.14182, places spin-$\tau$ in a
specific bidegree of a \emph{categorified} Hecke--Clifford trace --- which
need not descend to a per-basis-vector statement. The bidegree-detection
viewpoint is not refuted by the present probe.

\paragraph{(3) Different chain product.}
Crucially, the type-A definition of $\Om$ uses a $\hat\lambda$-cut: only the
factors $(T_{\ell+1}+1) \cdots (T_{n-1}+1)$ corresponding to the $1^{q}$ tail
enter. The probe in this paper used the \emph{full} chain $\Omc =
(\Tc_{n-1}+1) \cdots (\Tc_{1}+1)$, because strict partitions have no
$1^{q}$ tail. The spin analogue presumably needs a different cut --- a
staircase tail $(k, k-1, \dots, 1)$ within $\nu$? a $\hat\nu$-cut along a
distinguished horizontal? The present refutation is for the naive full
chain, and a tailored cut is the most natural follow-up.

\paragraph{(4) Other Frobenius-super algebras.}
The dream-2 universalisation conjecture targeted not just Hecke--Clifford
but the entire family $\mathcal{W}^{(1)}_{n}(A)$ for Frobenius
superalgebras $A$ (Aida--Kimura BKP, Moran higher-level wreath). The
non-negativity of contents (Lemma~\ref{lem:nonneg-contents}) is specific to
shifted-content combinatorics; other Frobenius-super families with
\emph{signed} content systems (e.g.\ BKP with charge $\delta \ne 0$) need
their own probes.

\section{Methodological note}
\label{sec:method}

This refutation joins the pattern documented in
\verb|topics/methodological-principles.md|:
\begin{itemize}
\item \emph{Principle 7 (test dream-named frameworks empirically).}
A categorical / super conjecture posed by a dream should be tested by a
one-day SymPy probe before any deep categorification is attempted. The
probe converges in $O(\text{hours})$ and either refutes the per-basis claim
cleanly or supplies positive data that warrants the deep read.
The Dai--Zhang cell-match refutation (2026-05-19) was the first canonical
application of this principle; the present Spin-Diagnostic-4 refutation
is the second.
\item \emph{Eighth shadow of the dichotomy program.}
Counting from 2026-05-14: Hu--Mathas Ext-positivity, Goertzen--Williamson
v=1, BGG-L 1423-avoidance, Fischer--Gangl Pfaffian, BPD/Grothendieck,
Sandvik definition-only, Dai--Zhang cell-match, and now Spin-Diagnostic-4.
The consistent pattern: deep-theory machinery that gestures at the
right phenomenon but \emph{descends to a different level} than the
per-basis-vector diagnostic.
\end{itemize}

\section{Acknowledgements and files}

Computation in SymPy with mpmath backend (60 digits, $q = 1.7$), script
\verb|~/projects/scratch/2026-05-20-spin-diagnostic-4.py|.
Chen--Wan seminormal specification from arXiv:2605.11778 (sections 2--4).
Spin-content conventions cross-checked against the
\verb|draft-2026-05-20-spin-d4-plan.md| design note. The dream-2 motivation
trace is recorded in \verb|memory/frobenius-super-universal-hub.md|.

\end{document}
